\chapter{Dataset e Sequenze MRI}
In questo capitolo viene presentato il dataset utilizzato nel progetto, in riferimento alle caratteristiche delle sequenze di risonanza magnetica impiegate nella pipeline di elaborazione.
La comprensione delle proprietà delle diverse modalità di acquisizione è infatti fondamentale per interpretare le informazioni contenute nelle immagini e per implementare correttamente gli algoritmi di segmentazione.

\section{Dataset BraTS (Medical Segmentation Decathlon)}
Per l'analisi è stato impiegato il dataset BraTS (\textit{Brain Tumor Segmentation Benchmark 2016/2017}), che rappresenta uno standard di riferimento nella comunità scientifica per la valutazione di algoritmi di segmentazione automatica in ambito medicale~\cite{menze2015}. Nello specifico, contiene volumi MRI multimodali annotati da medici esperti.

Il dataset è disponibile pubblicamente sulla piattaforma \textit{Medical Segmentation Decathlon} (MSD)\footnote{\url{http://medicaldecathlon.com/}}, una competizione progettata per valutare le prestazioni di diversi algoritmi di \textit{Image Processing} su molteplici task~\cite{antonelli_medical_2022}.

Per il lavoro proposto, è stata utilizzata la task \textit{Task01\_BrainTumor}, relativa alla segmentazione di gliomi cerebrali.

Dal punto di vista strutturale, il dataset è organizzato in tre sottoinsiemi principali:
\begin{itemize}
    \item \textbf{imagesTr}: contiene i volumi MRI dei pazienti utilizzati per il training;
    \item \textbf{labelsTr}: contiene le corrispondenti ground truth annotate manualmente dai medici;
    \item \textbf{imagesTs}: contiene i volumi di test, senza annotazioni.
\end{itemize}

Nel presente lavoro sono stati utilizzati esclusivamente i dati annotati (\textit{imagesTr} e \textit{labelsTr}), così da valutare le prestazioni della pipeline mediante confronto diretto con la ground truth. Il sottoinsieme di test non è stato considerato, poiché tipicamente destinato alla validazione di approcci di \textit{Machine Learning}.

Il dataset di training comprende 484 pazienti, ciascuno identificato da un file in formato \textit{.nii.gz} contenente un volume multimodale. Ogni volume è rappresentato come un tensore quadridimensionale, in cui le prime tre dimensioni corrispondono allo spazio tridimensionale dell'immagine, mentre la quarta dimensione rappresenta le diverse sequenze MRI acquisite.

Nella versione MSD, i volumi sono stati ricampionati al fine di ottenere una risoluzione spaziale con \textit{voxel} (pixel volumetrici) di $1 \times 1 \times 1$ mm ed è stato effettuato lo \textit{skull-stripping}, ovvero la rimozione del cranio dall'acquisizione. Inoltre, le diverse modalità MRI sono state normalizzate in intensità e posizionate nello stesso sistema di riferimento spaziale. A seguito di tali operazioni, tutti i volumi presentano dimensioni uniformi di $240 \times 240 \times 155$ voxel e un numero di \textit{slice} pari a 155 per tutti i pazienti.

Le annotazioni di riferimento (ground truth) sono fornite come maschere segmentate, in cui i voxel sono etichettati secondo diverse classi patologiche, tra cui \textit{edema peritumorale}, \textit{core tumorale non-enhancing} (necrotico) e \textit{tumore enhancing} (attivo), oltre allo sfondo. Nella pipeline, tali etichette sono state binarizzate per ottenere una segmentazione di tipo tumore/non-tumore, per semplificarne l'implementazione.

\section{Sequenze MRI}
Ogni modalità di risonanza magnetica presenta specifiche proprietà di contrasto, che evidenziano differenti caratteristiche dei tessuti e della struttura della lesione ~\cite{menze2015}.

\begin{itemize}
    \item \textbf{FLAIR} (\textit{Fluid Attenuated Inversion Recovery});
    \item \textbf{T1} (T1 pesata);
    \item \textbf{T1c} (T1 con mezzo di contrasto);
    \item \textbf{T2} (T2 pesata).
\end{itemize}

\subsection{T1-pesata (T1)}
Le immagini \textit{T1} pesate sono caratterizzate da un buon contrasto tra sostanza grigia, sostanza bianca e liquido cerebrospinale. In questa modalità, i tessuti appaiono con intensità differenti, consentendo una visualizzazione chiara della struttura anatomica globale.

Risulta utile per la localizzazione generale delle strutture cerebrali, ma è poco sensibile nella rilevazione delle regioni tumorali, in quanto molte lesioni presentano intensità simili ai tessuti circostanti. Infatti, in assenza di mezzo di contrasto, le aree tumorali risultano difficilmente distinguibili.

\subsection{T1 con mezzo di contrasto (T1c)}
Nella sequenza \textit{T1c}, la somministrazione di un agente di contrasto a base di gadolinio evidenzia le aree tumorali attive, caratterizzate da una permeabilità vascolare aumentata.

Questa proprietà permette al mezzo di contrasto di accumularsi nei tessuti, rendendo tali regioni iperintense rispetto a quelle sane circostanti. Difatti, questa sequenza è efficace nell’identificazione della porzione attiva della lesione, ma non fornisce informazioni sull’estensione totale del tumore.

\subsection{T2-pesata (T2)}
Le immagini \textit{T2} pesate sono sensibili al contenuto di acqua nei tessuti. Di conseguenza, le regioni contenenti liquidi, come il liquido cerebrospinale o gli edemi, appaiono iperintense.

Tale sequenza contribuisce nell'individuazione dell’estensione della lesione, ma è meno specifica rispetto alla sequenza \textit{T1} con mezzo di contrasto, in quanto non consente di distinguere in modo accurato le diverse componenti tumorali. Per questo motivo, le immagini \textit{T2} risultano utili se integrate con altre modalità MRI complementari.

\subsection{FLAIR}
La sequenza \textit{FLAIR} sopprime il segnale del liquido cerebrospinale, mantenendo alta la visibilità delle anomalie nei tessuti cerebrali.

Essa è utile per localizzare l’edema peritumorale, che appare come una regione iperintensa estesa attorno alla lesione. A differenza della sequenza \textit{T1} con mezzo di contrasto che evidenzia la componente tumorale attiva, la \textit{FLAIR} consente di osservare l’estensione complessiva del tumore, includendo anche le aree non attive.

\section{Fusione Multimodale}
Come descritto, l’integrazione di più modalità MRI permette di migliorare le prestazioni degli algoritmi di segmentazione. Tuttavia, ciascuna sequenza presenta vantaggi e limitazioni specifiche.

Pertanto, nella pipeline è stata adottata una strategia di combinazione pesata delle sequenze per ogni paziente. In particolare, è stata implementata una tecnica di fusione a livello di intensità dei pixel prima della fase di segmentazione, come presentato in alcuni lavori di ricerca~\cite{james2014}.

Tale approccio fornisce in ingresso alla pipeline un’immagine arricchita, che valorizza le informazioni più rilevanti e riduce le ambiguità legate alle singole modalità.